\scanpage{134}
\mspage{50}

\deleted{\unsure{ça|ça.} ce sont les exemples qui manquent. La clef d'une compréhension semble bien être dans les questions de \unsure{fonctorialité|fonctorisation} entre ensembles $\Phi\longleftrightarrow\Psi$. Il faudrait en faire une étude systématique, mais je recule devant la perspective d'une digression qui me prendrait peut-être encore des semaines !}

4. \inserted{\uline{Paratopos et}} \uline{L'opération} $M\boxtimes N$.

4.1. \inserted{\uline{Pseudo-topos.} \uline{Pb. du \unsure{produit|premier}}}

J'appelle \inserted{\uline{Catégorie}} \uline{pseudo-$\unsure{U|\mathcal U}$-topos}, ou simplement \uline{pseudo-topos}, une \inserted{$U$-}catégorie $M$ stable par $U$-limites \deleted{projectives}\inserted{inductives}, et qui admet une \uline{petite} sous-catégorie $M_0$ \deleted{stable par} \uline{strictement génératrice} (i.e. génératrice par \unsure{épimorphismes} stricts). Alors $M$ admet aussi \unsure{ses|des} $U$-limites projectives \inserted{(SGA 4 I 8.\unsure{12|11}.\unsure{7|3})}, \deleted{[Mais} l'exemple de (\unsure{Gr|An|Cn}), cf loc cit 8.12.9, montre que $M^{\mathrm{o}}$ n'est pas forcément, pour autant, un pseudo-$U$-topos, faute à $M$ d'admettre une petite \deleted{\unsure{famille}} sous-catégorie strictement \uline{co}génératrice. Mais le th. 2.6. donne des conditions sur $M$ (\unsure{surtout} du type \unsure{exactitude})
% archive sheet number: 133
